%
% 第5章：核心操作

\clearpage
\cleardoublepage

\chapter{神经网络核心操作}

\section{引言}

本章介绍现代神经网络的基本构建块：卷积、池化和归一化操作。理解这些操作对于设计和实现有效的神经网络架构至关重要。

\section{卷积操作}

卷积是卷积神经网络（CNN）的基石。

对于输入 $\mathbf{X} \in \mathbb{R}^{H \times W}$ 和卷积核 $\mathbf{K} \in \mathbb{R}^{K \times K}$：
\begin{equation}
y[i,j] = \sum_{m=0}^{K-1}\sum_{n=0}^{K-1} k[m,n] \cdot x[i+m, j+n]
\end{equation}

卷积参数：
\begin{itemize}
    \item \textbf{卷积核大小：} $K \times K$（常见：$3\times3$, $5\times5$, $7\times7$）
    \item \textbf{步长（$s$）：} 滑动窗口的步长
    \item \textbf{填充（$p$）：} 输入周围的零填充
    \item \textbf{膨胀（$d$）：} 卷积核元素之间的间距
\end{itemize}

输出大小：
\begin{equation}
H_{out} = \left\lfloor \frac{H_{in} + 2p - d(K-1) - 1}{s} + 1 \right\rfloor
\end{equation}

卷积优势：
\begin{itemize}
    \item \textbf{局部感受野}
    \item \textbf{参数共享}
    \item \textbf{平移等变性}
\end{itemize}

\section{池化操作}

池化减少空间维度并提供平移不变性。

\subsection{最大池化}

\begin{equation}
y[i,j] = \max_{m,n \in \mathcal{N}(i,j)} x[m,n]
\end{equation}

保留最强激活，提供局部平移不变性。

\subsection{平均池化}

\begin{equation}
y[i,j] = \frac{1}{K \cdot K}\sum_{m,n \in \mathcal{N}(i,j)} x[m,n]
\end{equation}

在邻域内取平均。

\section{归一化技术}

归一化通过控制层输入分布来稳定训练。

\subsection{批归一化（BatchNorm）}

在批次维度上归一化激活：
\begin{equation}
\hat{x}_i = \frac{x_i - \mu_B}{\sqrt{\sigma_B^2 + \epsilon}}, \quad y_i = \gamma \hat{x}_i + \beta
\end{equation}

属性：
\begin{itemize}
    \item \textbf{减少内部协变量偏移}
    \item \textbf{允许更高的学习率}
    \item \textbf{提供轻微正则化}
\end{itemize}

\subsection{层归一化（LayerNorm）}

在单个样本的特征维度上归一化：
\begin{equation}
\mu = \frac{1}{D}\sum_{i=1}^{D} x_i, \quad \sigma = \frac{1}{D}\sum_{i=1}^{D}(x_i - \mu)^2
\end{equation}

独立于批次大小，用于RNN和Transformer。

\begin{table}[h]
    \centering
    \caption{归一化方法比较}
    \begin{tabular}{L{2cm} L{2.5cm} L{2.5cm} L{3cm}}
        \toprule
        \textbf{方法} & \textbf{归一化维度} & \textbf{依赖批次} & \textbf{使用场景} \\
        \midrule
        BatchNorm & $(H, W, C)$ & 是 & CNNs \\
        LayerNorm & $(H, W, N)$ & 否 & Transformers \\
        InstanceNorm & $(H, W)$ & 否 & 风格迁移 \\
        GroupNorm & $(H, W, G)$ & 否 & 小批次CNNs \\
        \bottomrule
    \end{tabular}
\end{table}

\section{本章小结}

\begin{enumerate}
    \item 卷积通过共享权重提取局部特征
    \item 池化提供空间约简和平移不变性
    \item 归一化稳定训练动态
    \item 现代架构将这些操作组合成标准化块
    \item 操作选择取决于任务和数据特征
\end{enumerate}
